\chapter{Basis, Dimensi, dan Transformasi Linear}

Sebuah detektor anomali tidak membaca lalu lintas jaringan mentah. Ia meringkasnya lebih
dulu. Empat pencacah, misalnya koneksi gagal, paket keluar, lama sesi, dan banyaknya porta
unik, diringkas menjadi tiga fitur oleh peta linear
$F\colon\mathbb{R}^{4}\to\mathbb{R}^{3}$.
Ambang keputusan hanya melihat $F(\mathbf{x})$.

Pertanyaan yang menentukan: adakah $\boldsymbol{\delta}\neq\mathbf{0}$ dengan
$F(\mathbf{x}+\boldsymbol{\delta})=F(\mathbf{x})$? Karena $F$ linear, syarat itu setara
dengan
$F(\boldsymbol{\delta})=\mathbf{0}$. Jawabannya subruang berdimensi
$4-\operatorname{rank}F$; bila $F$ berpangkat tiga, ada satu garis penuh arah buta, dan
penyerang dapat menggeser keempat pencacah sepanjang garis itu tanpa menggeser satu pun
fitur.

Untuk menjawabnya kita perlu tiga hal yang belum diberikan Bab 6: daftar koordinat terhadap
basis, aturan memindahkan daftar itu ketika basis berganti, dan terjemahan peta linear
menjadi matriks sehingga inti dapat dihitung dengan eliminasi. Bab 6 sudah menetapkan bahwa
basis ada dan ukurannya tunggal; aksiomanya tidak diulang di sini.

\begin{untukanak}
Bayangkan sebuah lemari buku dengan banyak rak. Setiap buku punya tempat tetap. Kalau kita
menata ulang lemarinya, buku-bukunya tidak berubah, tetapi alamat setiap buku berubah. Untuk
mengatakan di mana sebuah buku berada, kita harus menyebut penataan yang sedang dipakai.
Kawan kita yang memakai penataan lain akan menyebut alamat yang berbeda untuk buku yang
sama,
dan kalian berdua sama-sama benar.

Ada juga rak khusus di belakang lemari. Buku yang dipindahkan ke sana tidak terlihat dari
depan. Selama buku hanya berpindah di dalam rak itu, daftar buku di depan tidak berubah sama
sekali.

Pernahkah kamu memindahkan sesuatu ke tempat yang tidak diperiksa siapa pun? Kita akan tahu
mengapa daftar alamat bisa berubah sementara isi
lemari tetap, mengapa rak tersembunyi itu selalu punya ukuran tertentu, dan mengapa ada
perubahan yang tidak menggeser ringkasan sama sekali.
\end{untukanak}

\section{Koordinat dan matriks transisi}
\tingkat{menengah}

Vektor di $\mathbb{R}^{3}$ sudah berupa daftar tiga angka; vektor di $P_{2}$ atau
$M_{2\times2}(\mathbb{R})$ tidak. Supaya keduanya dapat dihitung dengan cara yang sama, kita
perlu terjemahan yang mengubah setiap vektor menjadi daftar angka. Terjemahan itu bergantung
pada basis, dan di situlah kesulitan bagian ini bermula.

Ketunggalan sudah kita punya: lemma ketunggalan koordinat di Bab 6 menyatakan penulisan
vektor dalam suatu basis selalu tunggal. Karena itu definisi berikut sah tanpa memilih di
antara beberapa penulisan.

\kamus{koordinat}{coordinate vector}{Alamat sebuah buku di lemari: isinya tetap, alamatnya
bergantung cara menata rak.}

\begin{definisi}[Vektor koordinat]
Misalkan $B=(\mathbf{b}_{1},\dots,\mathbf{b}_{n})$ adalah basis terurut dari ruang vektor
$V$
atas $\mathbb{R}$. Untuk setiap $\mathbf{v}\in V$ ada skalar tunggal $a_{1},\dots,a_{n}$
dengan $\mathbf{v}=a_{1}\mathbf{b}_{1}+\dots+a_{n}\mathbf{b}_{n}$. Kolom
\[
  [\mathbf{v}]_{B}=(a_{1},\dots,a_{n})^{\top}
\]
disebut \emph{vektor koordinat} $\mathbf{v}$ terhadap $B$.\index{koordinat}
\end{definisi}

Kata ``terurut'' bukan hiasan. Bila dua vektor basis ditukar, dua entri koordinat ikut
bertukar sedangkan vektornya tetap. Satu himpunan basis belum menentukan daftar koordinat;
urutannya juga menentukan.

Satu peringatan: $[\mathbf{v}]_{B}$ bukan $\mathbf{v}$. Angka-angka itu hidup di
$\mathbb{R}^{n}$, sedangkan $\mathbf{v}$ tinggal di $V$. Untuk $M_{2\times2}(\mathbb{R})$
bedanya terasa: vektornya berisi empat angka, dan koordinatnya juga empat angka, tetapi
tersusun sebagai kolom yang bergantung pada basis.

Sifat penting pertama dari terjemahan ini: ia menghormati kedua operasi ruang vektor.

\begin{proposisi}
Untuk setiap basis $B$ dari ruang $V$ berdimensi $n$, peta
$\varphi_{B}\colon V\to\mathbb{R}^{n}$ dengan $\varphi_{B}(\mathbf{v})=[\mathbf{v}]_{B}$
bersifat linear, injektif, dan surjektif.
\end{proposisi}

\begin{proof}
Ambil $\mathbf{v}=\sum_{i}a_{i}\mathbf{b}_{i}$ dan $\mathbf{w}=\sum_{i}c_{i}\mathbf{b}_{i}$.
Maka $\mathbf{v}+\mathbf{w}=\sum_{i}(a_{i}+c_{i})\mathbf{b}_{i}$. Ketunggalan koordinat
memaksa $[\mathbf{v}+\mathbf{w}]_{B}=(a_{i}+c_{i})=[\mathbf{v}]_{B}+[\mathbf{w}]_{B}$.
Dengan cara yang sama, $[c\mathbf{v}]_{B}=c[\mathbf{v}]_{B}$.

Untuk injektivitas, misalkan $[\mathbf{v}]_{B}=\mathbf{0}$. Maka semua koefisiennya nol,
sehingga $\mathbf{v}=\sum_{i}0\cdot\mathbf{b}_{i}=\mathbf{0}$. Untuk surjektivitas, ambil
sembarang $(a_{1},\dots,a_{n})\in\mathbb{R}^{n}$ dan bentuk
$\mathbf{v}=\sum_{i}a_{i}\mathbf{b}_{i}$; koordinatnya tepat daftar itu.
\end{proof}

Peta $\varphi_{B}$ memindahkan pertanyaan tentang $V$ menjadi pertanyaan tentang kolom di
$\mathbb{R}^{n}$, lalu memindahkannya kembali; pemindahan itu bergantung pada $B$. Koordinat
ditulis sebagai kolom agar matriks bekerja dari kiri: koordinat masuk, dikalikan matriks,
koordinat keluar. Dengan koordinat baris, setiap rumus harus ditranspos.

\begin{bacakeras}
Dua suara berdiri di depan kelas.

``Berapa panjang meja ini?''
``Tiga.''
``Tiga apa?''
``Tiga penggaris. Atau tiga jengkal. Atau tiga langkah. Pilih dulu alat ukurnya.''

Begitu juga ruang vektor. Sebuah vektor belum punya alamat sampai kita memilih basis. Basis
itu
penggaris: ia memberi skala, memberi arah, dan memberi urutan. Ganti penggarisnya, ganti
angkanya. Benda yang diukur tidak ikut berubah.

Satu penggaris, satu angka. Dua penggaris, dua angka. Tanpa penggaris, tidak ada angka sama
sekali.

Hati-hati dengan penggaris yang jelek: dua garis skalanya hampir sejajar. Angka tetap
keluar,
tetapi angkanya gemetar, dan dua orang yang mengukur dengan sabar akan mendapat angka
berbeda.

Ulangi bersama: basis bukan bendanya, basis alat ukurnya. Yang berubah alamat, yang tetap
isi.
\end{bacakeras}

Sekarang dua basis dipakai pada ruang yang sama. Setiap vektor punya dua daftar koordinat,
dan kita ingin rumus yang mengubah daftar pertama menjadi daftar kedua. Matriks yang
melakukan tugas itu disusun dari koordinat vektor basis lama dalam basis baru.

\begin{definisi}[Matriks transisi]
Misalkan $B=(\mathbf{b}_{1},\dots,\mathbf{b}_{n})$ dan
$C=(\mathbf{c}_{1},\dots,\mathbf{c}_{n})$
dua basis dari $V$. \emph{Matriks transisi} dari $B$ ke $C$ adalah matriks $n\times n$
\[
P_{C\leftarrow
B}=\Bigl(\,[\mathbf{b}_{1}]_{C}\;\Big|\;\cdots\;\Big|\;[\mathbf{b}_{n}]_{C}\,\Bigr),
\]
yaitu matriks yang kolom ke-$j$-nya adalah $[\mathbf{b}_{j}]_{C}$.\index{transisi}
\end{definisi}

Yang perlu dibuktikan: mengalikan daftar koordinat lama dengan matriks ini menghasilkan
daftar
koordinat baru.

\begin{teorema}
Untuk setiap $\mathbf{v}\in V$ berlaku $[\mathbf{v}]_{C}=P_{C\leftarrow B}[\mathbf{v}]_{B}$.
\end{teorema}

\begin{proof}
Tulis $\mathbf{v}=\sum_{j}a_{j}\mathbf{b}_{j}$, sehingga
$[\mathbf{v}]_{B}=(a_{1},\dots,a_{n})^{\top}$.
Kelinearan $\varphi_{C}$ memberi
$[\mathbf{v}]_{C}=\sum_{j}a_{j}[\mathbf{b}_{j}]_{C}$. Ruas kanan itu tepat hasil kali
matriks
$P_{C\leftarrow B}$ dengan kolom $(a_{j})$, menurut definisi perkalian matriks dengan vektor.
\end{proof}

Akibat yang membuat matriks itu dapat dipakai dua arah.

\begin{akibat}
Matriks transisi selalu invertible, dan $P_{B\leftarrow C}=P_{C\leftarrow B}^{-1}$. Selain itu
$P_{C\leftarrow B}=P_{C\leftarrow D}\,P_{D\leftarrow B}$ untuk tiga basis $B,C,D$ dari $V$.
\end{akibat}

\begin{proof}
Teorema sebelumnya memberi $P_{B\leftarrow C}P_{C\leftarrow B}[\mathbf{v}]_{B}=[\mathbf{v}]_{B}$
untuk semua $\mathbf{v}$. Karena $\varphi_{B}$ surjektif, hasil kali itu sama dengan matriks
identitas. Argumen yang sama berlaku pada urutan sebaliknya, jadi kedua matriks saling
invers. Untuk aturan rantai, terapkan teorema dua kali:
$[\mathbf{v}]_{C}=P_{C\leftarrow D}[\mathbf{v}]_{D}=P_{C\leftarrow D}P_{D\leftarrow B}[\mathbf{v}]_{B}$,
lalu bandingkan dengan $[\mathbf{v}]_{C}=P_{C\leftarrow B}[\mathbf{v}]_{B}$ pada $\mathbf{v}=\mathbf{b}_{j}$.
\end{proof}

Perhatikan bahwa $P_{C\leftarrow B}$ dan $P_{B\leftarrow C}$ bukan transpos satu sama lain; keduanya saling invers, dan itu berbeda. Arah subskrip dibaca dari kanan ke kiri, seperti komposisi peta: ambil vektor basis ke-$j$ dari $B$, lalu tanyakan alamatnya menurut $C$. Bila $B$ basis standar, kolom-kolom $P_{B\leftarrow C}$ adalah vektor-vektor $C$ itu sendiri.

Ambil $B=(\mathbf{e}_{1},\mathbf{e}_{2},\mathbf{e}_{3})$
dan $C=(\mathbf{b}'_{1},\mathbf{b}'_{2},\mathbf{b}'_{3})$ dengan
\[
  \mathbf{b}'_{1}=(1,1,0),\qquad \mathbf{b}'_{2}=(1,0,1),\qquad \mathbf{b}'_{3}=(0,1,1).
\]
Karena $B$ standar, matriks $P_{B\leftarrow C}$ memuat ketiga vektor itu sebagai kolom:
\[
  P_{B\leftarrow C}=\begin{pmatrix}1&1&0\\ 1&0&1\\ 0&1&1\end{pmatrix},
  \qquad \det P_{B\leftarrow C}=-2\neq0 .
\]
Inversnya
\[
  P_{C\leftarrow B}=\frac{1}{2}\begin{pmatrix}1&1&-1\\ 1&-1&1\\ -1&1&1\end{pmatrix}.
\]
Untuk $\mathbf{v}=(3,8,11)$ kita peroleh
\[
  [\mathbf{v}]_{C}=P_{C\leftarrow B}[\mathbf{v}]_{B}
  =\frac{1}{2}\begin{pmatrix}3+8-11\\ 3-8+11\\ -3+8+11\end{pmatrix}
  =\begin{pmatrix}0\\ 3\\ 8\end{pmatrix}.
\]
Pemeriksaannya langsung: $3\mathbf{b}'_{2}+8\mathbf{b}'_{3}=(3,8,11)$. Sebaliknya
$\mathbf{w}=(1,0,0)$ berkoordinat $(1/2,\,1/2,\,-1/2)$; entri pecahan wajar, karena
koordinat
hanyalah alamat.

\section{Transformasi linear, inti, dan jangkauan}
\tingkat{dasar}

Bab 6 membahas ruang dan subruang, tetapi belum membahas peta di antaranya. Padahal
pekerjaan utama aljabar linear adalah memindahkan vektor: memutar, memproyeksikan,
menurunkan, meringkas.

Aturan yang menghormati penjumlahan dan perkalian skalar disebut transformasi linear:
menjumlahkan dulu lalu memetakan sama hasilnya dengan memetakan dulu lalu menjumlahkan.

\begin{definisi}[Transformasi linear]
Peta $T\colon V\to W$ antara dua ruang vektor atas $\mathbb{R}$ disebut
\emph{transformasi linear} bila $T(\mathbf{u}+\mathbf{v})=T(\mathbf{u})+T(\mathbf{v})$ dan
$T(c\mathbf{v})=cT(\mathbf{v})$ untuk semua $\mathbf{u},\mathbf{v}\in V$ dan $c\in\mathbb{R}$.
\end{definisi}

Kedua syarat setara dengan satu syarat,
$T(c\mathbf{u}+\mathbf{v})=cT(\mathbf{u})+T(\mathbf{v})$. Konsekuensi pertama,
$T(\mathbf{0})=\mathbf{0}$, dipakai untuk menolak kelinearan dengan cepat: peta
$T(\mathbf{x})=\mathbf{x}+\mathbf{a}$ dengan $\mathbf{a}\neq\mathbf{0}$ langsung gugur.

Sebaliknya $T(\mathbf{0})=\mathbf{0}$ tidak cukup. Peta $T(x,y)=(x^{2},y)$ memenuhinya,
tetapi
$T(2,0)=(4,0)$ sedangkan $2T(1,0)=(2,0)$. Kuadrat memutus kelinearan.\index{transformasi
linear}

Rotasi $90^{\circ}$, $R(x,y)=(-y,x)$, linear: $R(3,1)=(-1,3)$. Turunan $D\colon P_{3}\to P_{2}$, $D(p)=p'$, linear karena $(p+q)'=p'+q'$. Peta $T(x,y)=(x+y,\,2x+2y)$ juga linear, tetapi seluruh $\mathbb{R}^{2}$ jatuh ke satu garis; ia linear sekaligus merusak informasi.

Ambil
$A=\begin{pmatrix}2&1\\1&2\end{pmatrix}$ dan $T(\mathbf{x})=A\mathbf{x}$. Vektor $(1,0)$ pergi
ke $(2,1)$ dan $(0,1)$ ke $(1,2)$, sehingga persegi satuan menjadi jajar genjang bertitik
sudut
$(0,0)$, $(2,1)$, $(3,3)$, $(1,2)$ dengan luas $3$, sebab $|2\cdot2-1\cdot1|=3$.
Gambar~\ref{fig:kisi-07} memperlihatkan kisi sebelum dan sesudah peta: garis tetap garis, kisi
tetap kisi, tetapi jarak dan luas tidak lagi sama.

\begin{figure}[htbp]
\centering
\begin{tikzpicture}[scale=0.38,>=stealth,font=\scriptsize]
  \foreach \k in {0,1,2,3}{
    \draw[layergrey!45] (\k,0)--(\k,3);
    \draw[layergrey!45] (0,\k)--(3,\k);
    \draw[layerteal,thick] (2*\k,\k)--(2*\k+3,\k+6);
    \draw[layerteal,thick] (\k,2*\k)--(\k+6,2*\k+3);}
  \fill[accent!15] (0,0)--(1,0)--(1,1)--(0,1)--cycle;
  \draw[accent,thick] (0,0)--(1,0)--(1,1)--(0,1)--cycle;
  \fill[layerorange!25] (0,0)--(2,1)--(3,3)--(1,2)--cycle;
  \draw[layerorange,very thick] (0,0)--(2,1)--(3,3)--(1,2)--cycle;
  \draw[->,ink,thick] (-0.6,0)--(9.8,0) node[below right]{$x$};
  \draw[->,ink,thick] (0,-0.6)--(0,9.8) node[left]{$y$};
  \foreach \x/\y in {2/1,1/2,3/3}{
    \fill[ink] (\x,\y) circle (1.8pt);
    \node[ink,anchor=north west] at (\x+0.05,\y-0.05) {$(\x,\y)$};}
  \node[layergrey,anchor=north west] at (0,-1.0) {tipis: kisi asal};
  \node[layerteal,anchor=north west] at (4.0,-1.0) {tebal: kisi setelah dipetakan};
\end{tikzpicture}
\caption{Kisi satuan dan bayangannya di bawah $T(\mathbf{x})=A\mathbf{x}$ dengan
$A=\begin{pmatrix}2&1\\1&2\end{pmatrix}$. Satu petak berubah menjadi jajar genjang bertitik
sudut $(0,0)$, $(2,1)$, $(3,3)$, dan $(1,2)$, sehingga luasnya naik dari $1$ menjadi $3$.
Kisi lama digambar tipis, kisi baru digambar tebal; keduanya tetap kisi.}
\label{fig:kisi-07}
\end{figure}

\begin{tangan}
Ambil selembar kertas berpetak, pensil, dan penggaris. Gambar sumbu mendatar dan sumbu
tegak,
lalu tandai perpotongannya sebagai pusat. Gambar sebuah rumah kecil: persegi lima petak,
atap
segitiga, pintu satu petak.

Sekarang geser. Pindahkan setiap titik rumah itu tiga petak ke kanan dan dua petak ke atas.
Jangan menebak; hitung. Titik yang tadinya satu petak ke kanan dan satu petak ke atas
sekarang
berada empat petak ke kanan dan tiga petak ke atas. Gambar ulang rumah itu. Bentuk dan
ukurannya
sama, hanya tempatnya berpindah.

Sekarang putar. Putar kertas seperempat putaran searah jarum jam, lalu gambar ulang rumah
yang
sama. Panjang setiap sisi tetap, arah setiap sisi berubah.

Yang harus terlihat: pada geseran, semua arah bertahan tetapi pusat ikut berpindah; pada
putaran, pusat bertahan dan semua arah berubah. Karena geseran memindahkan pusat, geseran
bukan
transformasi linear. Putaran termasuk.
\end{tangan}

Kabar baiknya: untuk mengenal peta linear, cukup mengenal nilainya pada vektor basis.

\begin{proposisi}
Misalkan $B=(\mathbf{b}_{1},\dots,\mathbf{b}_{n})$ basis $V$ dan
$\mathbf{w}_{1},\dots,\mathbf{w}_{n}\in W$ sebarang. Ada tepat satu peta linear
$T\colon V\to W$ dengan $T(\mathbf{b}_{j})=\mathbf{w}_{j}$ untuk setiap $j$.
\end{proposisi}

\begin{proof}
Untuk keberadaan, definisikan
$T\bigl(\sum_{j}a_{j}\mathbf{b}_{j}\bigr)=\sum_{j}a_{j}\mathbf{w}_{j}$.
Definisi ini tidak bergantung pada penulisan, karena koordinat terhadap basis tunggal. Untuk
kelinearan, pakai kelinearan koordinat dari proposisi di bagian sebelumnya:
$T(\mathbf{u}+\mathbf{v})=\sum_{j}(a_{j}+c_{j})\mathbf{w}_{j}=T(\mathbf{u})+T(\mathbf{v})$,
dan serupa untuk perkalian skalar.

Untuk ketunggalan, misalkan $S$ linear dan $S(\mathbf{b}_{j})=\mathbf{w}_{j}$. Untuk
$\mathbf{v}=\sum_{j}a_{j}\mathbf{b}_{j}$ berlaku
$S(\mathbf{v})=\sum_{j}a_{j}S(\mathbf{b}_{j})=\sum_{j}a_{j}\mathbf{w}_{j}=T(\mathbf{v})$.
\end{proof}

Akibat praktisnya: $n$ vektor di $W$ menentukan satu peta linear dari $V$ berdimensi $n$.
Bila $W=\mathbb{R}^{m}$, data itu tidak lain matriks $m\times n$.

Setiap peta linear membawa dua himpunan yang layak diperiksa: vektor yang dipetakan ke nol,
dan nilai yang benar-benar tercapai.

\kamus{inti}{kernel atau null space}{Daftar perubahan yang sama sekali tidak menggeser
ringkasan.}

\begin{definisi}[Inti, jangkauan, rank, dan nullity]
Untuk $T\colon V\to W$ linear, \emph{inti} adalah
$\ker T=\{\mathbf{v}\in V: T(\mathbf{v})=\mathbf{0}\}$ dan \emph{jangkauan} adalah
$\operatorname{im}T=\{T(\mathbf{v}):\mathbf{v}\in V\}$. Bila $V$ berdimensi hingga,
$\operatorname{nullity}(T)=\dim\ker T$ dan $\operatorname{rank}(T)=\dim\operatorname{im}T$.
\end{definisi}

Keduanya subruang, bukan sekadar himpunan.\index{inti}\index{jangkauan}

\begin{lemma}
$\ker T$ subruang $V$ dan $\operatorname{im}T$ subruang $W$.
\end{lemma}

\begin{proof}
Keduanya memuat $\mathbf{0}$ karena $T(\mathbf{0})=\mathbf{0}$. Bila
$\mathbf{u},\mathbf{v}\in\ker T$
dan $c\in\mathbb{R}$, maka $T(\mathbf{u}+\mathbf{v})=\mathbf{0}$ dan
$T(c\mathbf{u})=\mathbf{0}$,
jadi $\ker T$ tertutup. Bila $\mathbf{u}=T(\mathbf{x})$ dan $\mathbf{v}=T(\mathbf{y})$, maka
$\mathbf{u}+\mathbf{v}=T(\mathbf{x}+\mathbf{y})$ dan $c\mathbf{u}=T(c\mathbf{x})$ anggota
$\operatorname{im}T$.
\end{proof}

Inti mengukur kegagalan injektivitas, dan hanya itu yang perlu diperiksa.

\begin{proposisi}
$T$ injektif jika dan hanya jika $\ker T=\{\mathbf{0}\}$.
\end{proposisi}

\begin{proof}
Bila $T$ injektif, hanya $\mathbf{0}$ yang punya peta $\mathbf{0}$. Sebaliknya andaikan
$\ker T=\{\mathbf{0}\}$ dan $T(\mathbf{u})=T(\mathbf{v})$. Maka $T(\mathbf{u}-\mathbf{v})=\mathbf{0}$,
jadi $\mathbf{u}-\mathbf{v}\in\ker T$, sehingga $\mathbf{u}=\mathbf{v}$.
\end{proof}

Rank menghitung arah yang bertahan, nullity menghitung arah yang lenyap. Keduanya membagi
dimensi ruang asal.

\kamus{rank}{rank}{Banyaknya arah bebas yang masih bertahan setelah dipetakan.}

Contoh kecil memperlihatkan kedua nasib itu. Ambil
$B=\begin{pmatrix}1&2\\2&4\end{pmatrix}$. Maka $B\mathbf{e}_{1}=(1,2)$ dan
$B\mathbf{e}_{2}=(2,4)$, dua vektor segaris, jadi $\operatorname{rank}(B)=1$. Arah $(1,2)$
bertahan dan memanjang: $B(1,2)=(5,10)$. Arah $(2,-1)$ lenyap: $B(2,-1)=(0,0)$.
Gambar~\ref{fig:arah} menggambar keduanya, dan rank-nullity menutup hitungan: $2=1+1$.

\begin{figure}[htbp]
\centering
\begin{tikzpicture}[scale=0.60,>=stealth,font=\scriptsize]
  \draw[layerjade!60,thin] (-0.7,-1.4)--(2.8,5.6);
  \node[layerjade,anchor=west,rotate=63] at (2.0,4.3) {garis arah $(1,2)$};
  \draw[->,ink,thick] (-1.0,0)--(3.3,0) node[below right]{$x$};
  \draw[->,ink,thick] (0,-2.0)--(0,5.0) node[left]{$y$};
  \draw[->,layergrey,thick] (0,0)--(1,0) node[below]{$\mathbf{e}_{1}$};
  \draw[->,layergrey,thick] (0,0)--(0,1) node[left]{$\mathbf{e}_{2}$};
  \draw[->,layerorange,very thick] (0,0)--(1,2) node[above left]{$B\mathbf{e}_{1}=(1,2)$};
  \draw[->,layerbrown,very thick] (0,0)--(2,4) node[right]{$B\mathbf{e}_{2}=(2,4)$};
  \draw[->,layerplum,thick,dashed] (0,0)--(1.8,-0.9) node[below right]{$(2,-1)$};
  \node[layerplum,anchor=north] at (1.5,-1.6) {runtuh ke $\mathbf{0}$};
  \fill[ink] (0,0) circle (1.6pt) node[below left,ink]{$\mathbf{0}$};
\end{tikzpicture}
\caption{Dua arah di bawah $B=\begin{pmatrix}1&2\\2&4\end{pmatrix}$. Arah $(1,2)$ bertahan dan
memanjang menjadi $(5,10)$; arah $(2,-1)$ runtuh ke $\mathbf{0}$. Bayangan kedua vektor basis,
$B\mathbf{e}_{1}=(1,2)$ dan $B\mathbf{e}_{2}=(2,4)$, terletak pada satu garis, jadi
$\operatorname{rank}(B)=1$ dan $\operatorname{nullity}(B)=1$.}
\label{fig:arah}
\end{figure}

Sekarang teorema pusat bab ini. Isinya: tidak ada dimensi yang hilang di tengah jalan.

\begin{teorema}[Rank-nullity]
Bila $V$ berdimensi hingga dan $T\colon V\to W$ linear, maka
\[
  \dim V=\dim\ker T+\dim\operatorname{im}T .
\]
\end{teorema}

\begin{proof}
Misalkan $\dim V=n$ dan $\dim\ker T=k$. Ambil basis
$(\mathbf{u}_{1},\dots,\mathbf{u}_{k})$ dari $\ker T$. Menurut akibat perluasan basis di Bab 6,
basis itu dapat dilengkapi menjadi basis $V$ dengan menambahkan $n-k$ vektor
$\mathbf{w}_{1},\dots,\mathbf{w}_{n-k}$.

Pertama, himpunan $\{T(\mathbf{w}_{1}),\dots,T(\mathbf{w}_{n-k})\}$ merentang
$\operatorname{im}T$. Untuk sembarang $T(\mathbf{v})$, tulis
$\mathbf{v}=\sum_{i}\alpha_{i}\mathbf{u}_{i}+\sum_{j}\beta_{j}\mathbf{w}_{j}$. Suku-suku
$\mathbf{u}_{i}$ lenyap karena $T(\mathbf{u}_{i})=\mathbf{0}$, sehingga
$T(\mathbf{v})=\sum_{j}\beta_{j}T(\mathbf{w}_{j})$.

Kedua, himpunan itu bebas linear. Andaikan $\sum_{j}\beta_{j}T(\mathbf{w}_{j})=\mathbf{0}$.
Maka $T\bigl(\sum_{j}\beta_{j}\mathbf{w}_{j}\bigr)=\mathbf{0}$, jadi
$\sum_{j}\beta_{j}\mathbf{w}_{j}\in\ker T$, dan ada skalar $\alpha_{i}$ dengan
$\sum_{j}\beta_{j}\mathbf{w}_{j}=\sum_{i}\alpha_{i}\mathbf{u}_{i}$. Pindahkan semua suku ke
satu ruas:
\[
  \sum_{j}\beta_{j}\mathbf{w}_{j}-\sum_{i}\alpha_{i}\mathbf{u}_{i}=\mathbf{0}.
\]
Relasi itu berada di dalam basis $V$, sehingga semua koefisiennya nol; khususnya setiap
$\beta_{j}=0$. Jadi $\{T(\mathbf{w}_{j})\}$ basis $\operatorname{im}T$, dan
$\operatorname{rank}(T)=n-k$.
\end{proof}

Gambar~\ref{fig:batang} menyusun teorema itu sebagai dua batang. Untuk ruang berdimensi
$4$ panjang totalnya selalu $4$: identitas memberi $4+0$, proyeksi yang membuang satu
koordinat
memberi $3+1$, dan peta fitur $G$ nanti memberi $2+2$. Yang bergeser hanya letak pemisahnya.

\begin{figure}[htbp]
\centering
\begin{tikzpicture}[x=1cm,y=1cm,font=\scriptsize,>=stealth]
  \fill[layerjade!45] (0,2.4) rectangle (4,3.1);
  \draw[layergrey] (0,2.4) rectangle (4,3.1);
  \node at (2,2.75) {rank $4$};
  \node[layerjade,anchor=east] at (-0.12,2.75) {identitas};
  \node[layerjade,anchor=west] at (4.12,2.75) {nullity $0$};
  \fill[layerteal!45] (0,1.2) rectangle (3,1.9);
  \fill[layerorange!45] (3,1.2) rectangle (4,1.9);
  \draw[layergrey] (0,1.2) rectangle (4,1.9);
  \draw[layergrey] (3,1.2)--(3,1.9);
  \node at (1.5,1.55) {rank $3$};
  \node at (3.5,1.55) {$1$};
  \node[layerteal,anchor=east] at (-0.12,1.55) {proyeksi};
  \node[layerorange,anchor=west] at (4.12,1.55) {nullity $1$};
  \fill[layerteal!45] (0,0) rectangle (2,0.7);
  \fill[layerorange!45] (2,0) rectangle (4,0.7);
  \draw[layergrey] (0,0) rectangle (4,0.7);
  \draw[layergrey] (2,0)--(2,0.7);
  \node at (1,0.35) {rank $2$};
  \node at (3,0.35) {$2$};
  \node[layerteal,anchor=east] at (-0.12,0.35) {peta fitur $G$};
  \node[layerorange,anchor=west] at (4.12,0.35) {nullity $2$};
  \draw[layergrey,dashed] (4,-0.25)--(4,3.4);
  \node[layergrey,anchor=south] at (4,3.4) {$\dim V=4$};
  \fill[layerteal!45] (0,-0.95) rectangle (0.35,-0.6);
  \draw[layergrey] (0,-0.95) rectangle (0.35,-0.6);
  \node[anchor=west] at (0.45,-0.78) {arah yang bertahan};
  \fill[layerorange!45] (3.3,-0.95) rectangle (3.65,-0.6);
  \draw[layergrey] (3.3,-0.95) rectangle (3.65,-0.6);
  \node[anchor=west] at (3.75,-0.78) {arah yang lenyap};
\end{tikzpicture}
\caption{Rank-nullity sebagai dua batang yang jumlahnya tetap. Untuk $\dim V=4$: peta identitas
memberi $4+0$, proyeksi yang membuang satu koordinat memberi $3+1$, dan peta fitur $G$ memberi
$2+2$. Hanya letak pemisahnya yang bergeser; panjang totalnya tetap $4$.}
\label{fig:batang}
\end{figure}

Dua akibat yang paling sering dipakai.

\begin{akibat}
Misalkan $\dim V=\dim W$ berhingga dan $T\colon V\to W$ linear. Maka $T$ injektif jika dan
hanya jika $T$ surjektif. Selain itu
$\operatorname{rank}(T)\le\min(\dim V,\dim W)$.
\end{akibat}

\begin{proof}
Menurut rank-nullity, $\dim\ker T=n-\operatorname{rank}T$ dengan $n=\dim V$. Injektif tepat
ketika $\dim\ker T=0$, yaitu $\operatorname{rank}T=n=\dim W$, yang berarti
$\operatorname{im}T=W$ karena subruang berdimensi penuh memenuhi seluruh ruang. Implikasi
sebaliknya diperoleh dengan membalik urutan argumen. Batas atas rank jelas dari
$\operatorname{im}T\subseteq W$ dan dari rank-nullity.
\end{proof}

Sebagai gambaran angka, ambil $T\colon\mathbb{R}^{3}\to\mathbb{R}^{2}$ dengan
$T(x,y,z)=(x-y,\,y-z)$. Karena $T(t,t,t)=(0,0)$, intinya garis $\{(t,t,t)\}$ dan nullity
$1$.
Bayangan $T(1,0,0)=(1,0)$ dan $T(0,0,1)=(0,-1)$ merentang $\mathbb{R}^{2}$, jadi rank $2$;
rank-nullity: $3=1+2$.

Bagi pemelajaran mesin, akibat itu keras. Peta fitur $\mathbb{R}^{n}\to\mathbb{R}^{m}$
dengan $n>m$ tidak mungkin injektif, dan intinya berdimensi paling sedikit $n-m$. Setiap
pemampatan dimensi menyimpan arah buta; pertanyaannya bukan apakah arah itu ada, melainkan
apakah penyerang dapat menemukannya.

Prosedur baku untuk basis inti: bawa matriks ke bentuk eselon tereduksi, tandai kolom
tanpa pivot sebagai variabel bebas, lalu gilir setiap variabel bebas dengan nilai satu dan
variabel bebas lain nol. Variabel pivot dibaca dari baris-barisnya. Hasilnya satu vektor per
kolom tanpa pivot, bebas linear karena masing-masing punya posisi satu yang tidak dimiliki
yang
lain.

Untuk jangkauan, prosedurnya berbeda dan sering tertukar dengan prosedur inti di atas. Empat
langkah berikut cukup dihafal.

\begin{resep}
\begin{enumerate}
  \item Susun matriks $A$, lalu bawa ke bentuk eselon baris tereduksi dengan eliminasi.
  \item Tandai kolom yang memuat pivot. Kolom tanpa pivot bukan calon anggota basis.
  \item Kembali ke matriks asli $A$, bukan ke bentuk eselonnya. Ambil kolom asli yang posisinya
        sama dengan kolom pivot.
  \item Periksa: banyaknya vektor terpilih sama dengan banyaknya baris tak nol, dan vektor itu
        bebas linear.
\end{enumerate}
Peringatan urutan: langkah tiga tidak boleh mengambil kolom dari bentuk eselon. Operasi
baris
mengubah kolom, sedangkan ruang kolom tidak berubah. Mengambil kolom dari matriks eselon
memberi jawaban salah walaupun banyaknya benar.
\end{resep}

Inti juga menjawab pertanyaan yang lebih umum: bila satu nilai sudah tercapai, semua
prapetanya tersusun rapi. Di sini $T^{-1}(\mathbf{w})$ berarti himpunan prapeta, bukan peta
invers.

\begin{proposisi}
Misalkan $T\colon V\to W$ linear dan $T(\mathbf{v}_{0})=\mathbf{w}$. Maka
$T^{-1}(\mathbf{w})=\{\mathbf{v}_{0}+\mathbf{k}:\mathbf{k}\in\ker T\}$.
\end{proposisi}

\begin{proof}
Bila $\mathbf{k}\in\ker T$, maka $T(\mathbf{v}_{0}+\mathbf{k})=\mathbf{w}+\mathbf{0}=\mathbf{w}$,
jadi setiap $\mathbf{v}_{0}+\mathbf{k}$ termasuk prapeta $\mathbf{w}$. Sebaliknya, bila
$T(\mathbf{v})=\mathbf{w}$, maka $T(\mathbf{v}-\mathbf{v}_{0})=\mathbf{0}$, sehingga
$\mathbf{v}-\mathbf{v}_{0}\in\ker T$ dan $\mathbf{v}=\mathbf{v}_{0}+(\mathbf{v}-\mathbf{v}_{0})$.
\end{proof}

Himpunan masukan yang menghasilkan fitur sama adalah salinan inti yang digeser: bila intinya
berdimensi satu, keluarganya satu garis, yaitu arah buta pada pembuka bab; bila berdimensi
$k$,
penyerang punya $k$ arah.

\section{Matriks representasi dan komposisi}
\tingkat{menengah}

Sekarang dua alat sebelumnya digabung: koordinat mengubah vektor menjadi kolom, dan peta
linear ditentukan nilainya pada basis. Gabungan keduanya menghasilkan matriks, tempat
seluruh
perhitungan dilakukan.

Matriks representasi menyimpan bayangan setiap vektor basis sebagai satu kolom, ditulis
dalam koordinat basis tujuan.

\begin{definisi}[Matriks representasi]
Misalkan $B=(\mathbf{b}_{1},\dots,\mathbf{b}_{n})$ basis $V$,
$C=(\mathbf{c}_{1},\dots,\mathbf{c}_{m})$ basis $W$, dan $T\colon V\to W$ linear.
\emph{Matriks representasi} $T$ terhadap $B$ dan $C$ adalah matriks $m\times n$
\[
[T]_{C\leftarrow
B}=\Bigl(\,[T(\mathbf{b}_{1})]_{C}\;\Big|\;\cdots\;\Big|\;[T(\mathbf{b}_{n})]_{C}\,\Bigr).
\]
\end{definisi}

Kolom-kolom itu diperoleh dengan menyelesaikan sistem, bukan menebak: untuk setiap $j$,
selesaikan $\sum_{i}x_{i}\mathbf{c}_{i}=T(\mathbf{b}_{j})$, lalu tulis solusinya sebagai
kolom
ke-$j$. Bila $C$ basis standar, penyelesaian itu tidak perlu dikerjakan, karena koordinat
terhadap basis standar adalah vektornya sendiri.

Klaim berikut menjelaskan mengapa matriks itu berguna: ia bekerja dengan cara yang
diharapkan.

\begin{teorema}
Untuk setiap $\mathbf{v}\in V$ berlaku
$[T(\mathbf{v})]_{C}=[T]_{C\leftarrow B}\,[\mathbf{v}]_{B}$.
\end{teorema}

\begin{proof}
Tulis $\mathbf{v}=\sum_{j}a_{j}\mathbf{b}_{j}$, jadi $[\mathbf{v}]_{B}=(a_{j})$. Kelinearan
$T$ memberi $T(\mathbf{v})=\sum_{j}a_{j}T(\mathbf{b}_{j})$. Pengambilan koordinat juga
linear,
sehingga $[T(\mathbf{v})]_{C}=\sum_{j}a_{j}[T(\mathbf{b}_{j})]_{C}$, yaitu tepat hasil kali
matriks dengan kolom $(a_{j})$.
\end{proof}

Dua akibat langsung. Peta linear dan matriksnya saling menentukan, karena nilai pada
basis sudah cukup. Dan rank tidak berubah dalam penerjemahan itu.

\begin{proposisi}
$\operatorname{rank}(T)$ sama dengan rank matriks $[T]_{C\leftarrow B}$, dan $\ker T$
berkorespondensi dengan ruang nol matriks itu. Rank matriks yang dimaksud adalah dimensi
ruang
kolomnya, sebagaimana sudah dipakai di Bab 6.
\end{proposisi}

\begin{proof}
Kolom ke-$j$ matriks adalah koordinat $T(\mathbf{b}_{j})$, dan
$\{T(\mathbf{b}_{j})\}$ merentang $\operatorname{im}T$. Pengambilan koordinat mempertahankan
kombinasi linear, jadi dimensi jangkauan sama dengan dimensi ruang kolom. Teorema sebelumnya
menunjukkan $T(\mathbf{v})=\mathbf{0}$ tepat ketika $[T(\mathbf{v})]_{C}=[T]_{C\leftarrow B}[\mathbf{v}]_{B}=\mathbf{0}$,
yaitu ketika $[\mathbf{v}]_{B}$ berada di ruang nol matriks itu.
\end{proof}

Komposisi peta berpadanan dengan perkalian matriks. Itu bukan kebetulan, melainkan
alasan perkalian matriks didefinisikan seperti sekarang.

\begin{teorema}
Misalkan $T\colon V\to W$ dan $S\colon W\to U$ linear, dengan basis $B,C,E$ pada ketiga ruang.
Maka $[S\circ T]_{E\leftarrow B}=[S]_{E\leftarrow C}\,[T]_{C\leftarrow B}$.
\end{teorema}

\begin{proof}
Untuk setiap $\mathbf{v}\in V$,
\[
  [(S\circ T)(\mathbf{v})]_{E}
  =[S]_{E\leftarrow C}[T(\mathbf{v})]_{C}
  =[S]_{E\leftarrow C}[T]_{C\leftarrow B}[\mathbf{v}]_{B}.
\]
Jadi matriks $[S]_{E\leftarrow C}[T]_{C\leftarrow B}$ memenuhi definisi matriks representasi
untuk $S\circ T$. Matriks representasi tunggal: bila dua matriks memberi hasil sama pada
semua $\mathbf{v}$, hasilnya sama khususnya pada $\mathbf{v}=\mathbf{b}_{j}$, sehingga
kolom-kolomnya sama.
\end{proof}

Komposisi juga memberi uji invertibilitas lewat matriks.

\begin{proposisi}
Misalkan $B$ dan $C$ basis dari ruang $V$ berdimensi $n$, dan $T\colon V\to V$ linear. Maka
$T$ invertible jika dan hanya jika $[T]_{C\leftarrow B}$ invertible, dan dalam hal itu
$[T^{-1}]_{B\leftarrow C}=([T]_{C\leftarrow B})^{-1}$.
\end{proposisi}

\begin{proof}
Bila $T$ invertible, maka $T^{-1}\circ T=I$ dan $T\circ T^{-1}=I$. Teorema komposisi memberi
$[T^{-1}]_{B\leftarrow C}[T]_{C\leftarrow B}=I$ dan
$[T]_{C\leftarrow B}[T^{-1}]_{B\leftarrow C}=I$, jadi kedua matriks saling invers.

Sebaliknya, misalkan $A=[T]_{C\leftarrow B}$ invertible. Bila $T(\mathbf{v})=\mathbf{0}$, maka
$A[\mathbf{v}]_{B}=\mathbf{0}$, sehingga $[\mathbf{v}]_{B}=\mathbf{0}$ dan
$\mathbf{v}=\mathbf{0}$.
Jadi $T$ injektif. Karena ruang asal dan ruang tujuan sama-sama berdimensi $n$, akibat
rank-nullity membuat $T$ surjektif, sehingga $T$ invertible. Persamaan untuk matriks
$T^{-1}$
lalu diikuti dari bagian pertama.
\end{proof}

Dengan begitu pertanyaan tentang invers operator menjadi pertanyaan tentang invers
matriks; jawaban tuntasnya ada di Bab 8.

Kasus khusus yang paling sering dipakai adalah operator pada satu ruang, $T\colon V\to V$,
dengan dua basis berbeda. Di sini kedua matriks terhubung oleh matriks transisi.

\begin{teorema}[Perubahan basis]
Misalkan $T\colon V\to V$ linear, $B$ dan $B'$ dua basis $V$, dan $P=P_{B\leftarrow B'}$.
Maka
\[
  [T]_{B'}=P^{-1}[T]_{B}P .
\]
\end{teorema}

\begin{proof}
Tulis $T=I\circ T\circ I$ dengan $I$ peta identitas. Menurut teorema komposisi,
\[
  [T]_{B'\leftarrow B'}=[I]_{B'\leftarrow B}\,[T]_{B\leftarrow B}\,[I]_{B\leftarrow B'}.
\]
Matriks identitas terhadap dua basis adalah matriks transisi, jadi
$[I]_{B'\leftarrow B}=P^{-1}$ dan $[I]_{B\leftarrow B'}=P$.
\end{proof}\index{representasi}

Dua matriks persegi yang terhubung seperti itu mendeskripsikan operator yang sama dari
sudut pandang berbeda, jadi keduanya diberi nama sendiri.

\begin{definisi}[Matriks serupa]
Matriks $A$ dan $B$ berukuran $n\times n$ disebut \emph{serupa} bila ada matriks invertible
$P$ dengan $B=P^{-1}AP$.\index{similaritas}
\end{definisi}

Similaritas karena itu nama lain untuk peta yang sama dilihat dari basis berbeda.

\begin{proposisi}
Dua matriks $n\times n$ serupa jika dan hanya jika keduanya merupakan matriks representasi
satu operator linear yang sama terhadap dua basis.
\end{proposisi}

\begin{proof}
Bila $B=P^{-1}AP$, ambil basis standar $B_{0}$ dan basis $B'$ yang kolom-kolomnya adalah
kolom-kolom $P$. Maka $P=P_{B_{0}\leftarrow B'}$, dan teorema perubahan basis memberi
$B=[T]_{B'}$ untuk operator $T(\mathbf{x})=A\mathbf{x}$. Sebaliknya, teorema yang sama
memberi
hubungan itu untuk setiap pasangan basis.
\end{proof}

Rank termasuk besaran yang tidak berubah: rank $[T]_{B'}$ sama dengan rank $T$, dan
begitu pula rank $[T]_{B}$. Determinan dan trace juga invarian, tetapi keduanya baru
didefinisikan di Bab 8 dan dibahas di Bab 9.

Contoh kecil untuk operator turunan: $D\colon P_{2}\to P_{2}$, $D(p)=p'$, basis
$B=(1,x,x^{2})$. Karena $D(1)=0$, $D(x)=1$, dan $D(x^{2})=2x$,
\[
  [D]_{B}=\begin{pmatrix}0&1&0\\ 0&0&2\\ 0&0&0\end{pmatrix}.
\]
Untuk $p=x^{2}+3x$ koordinatnya $(0,3,1)^{\top}$, dan matriks memberi $(3,\,2,\,0)^{\top}$,
yaitu $3+2x$, tepat turunan $p$.

\section{Isomorfisme, invarian dimensi, dan ruang peta linear}
\tingkat{lanjut}

Dua ruang berbeda dapat berperilaku sama. $P_{2}$ dan $\mathbb{R}^{3}$ berisi objek
berbeda jenis, tetapi keduanya punya tiga arah bebas dan aturan penjumlahan yang sama
bentuknya.

Korespondensi yang mempertahankan seluruh struktur itu disebut isomorfisme: bijektif,
linear, dan karena itu mempertahankan basis, dimensi, serta inti.

\kamus{isomorfisme}{isomorphism}{Dua denah rumah yang gambarnya berbeda tetapi jumlah dan
hubungan ruangannya sama.}

\begin{definisi}[Isomorfisme]
Peta linear $T\colon V\to W$ disebut \emph{isomorfisme} bila $T$ bijektif. Dua ruang disebut
\emph{isomorf} bila ada isomorfisme antara keduanya, ditulis $V\cong W$.
\end{definisi}

Satu lemma sebelum teorema utamanya: isomorfisme memindahkan basis menjadi basis.

\begin{lemma}
Bila $T\colon V\to W$ isomorfisme dan $B$ basis $V$, maka $T(B)$ basis $W$.
\end{lemma}

\begin{proof}
Misalkan $B=(\mathbf{b}_{1},\dots,\mathbf{b}_{n})$. Untuk kebebasan linear, andaikan
$\sum_{j}c_{j}T(\mathbf{b}_{j})=\mathbf{0}$. Karena $T$ linear, vektor
$\sum_{j}c_{j}\mathbf{b}_{j}$ dipetakan ke $\mathbf{0}$, jadi vektor itu ada di inti.
Intinya
hanya $\{\mathbf{0}\}$ karena $T$ injektif, sehingga semua $c_{j}=0$.
Untuk merentang, ambil $\mathbf{w}\in W$. Ada $\mathbf{v}\in V$ dengan $T(\mathbf{v})=\mathbf{w}$
karena $T$ surjektif, dan $\mathbf{v}=\sum_{j}a_{j}\mathbf{b}_{j}$ karena $B$ merentang.
Maka
$\mathbf{w}=\sum_{j}a_{j}T(\mathbf{b}_{j})$.
\end{proof}

Dengan lemma itu, kesamaan dimensi ternyata satu-satunya penghalang.

\begin{teorema}
Dua ruang berdimensi hingga atas $\mathbb{R}$ isomorf jika dan hanya jika dimensinya sama.
\end{teorema}

\begin{proof}
Bila $V\cong W$ lewat $T$, lemma sebelumnya mengubah basis $V$ menjadi basis $W$ dengan
banyaknya anggota sama, jadi $\dim V=\dim W$.

Sebaliknya, misalkan $\dim V=\dim W=n$. Pilih basis $B$ pada $V$ dan $C$ pada $W$. Peta
koordinat $\varphi_{B}\colon V\to\mathbb{R}^{n}$ dan $\varphi_{C}\colon W\to\mathbb{R}^{n}$
bijektif dan linear. Maka $\varphi_{C}^{-1}\circ\varphi_{B}$ adalah peta linear bijektif
dari
$V$ ke $W$, karena komposisi peta linear bijektif juga bijektif dan linear.
\end{proof}

Kesimpulan itu langsung terasa pada ruang yang kita kenal.

\begin{akibat}
Setiap ruang berdimensi $n$ isomorf dengan $\mathbb{R}^{n}$. Secara khusus
$P_{n-1}\cong\mathbb{R}^{n}$ dan $M_{2\times2}(\mathbb{R})\cong\mathbb{R}^{4}$.
\end{akibat}

Isomorfisme mempertahankan segala sesuatu yang didefinisikan dari operasi: kebebasan
linear, rentang, inti, jangkauan, rank, dan nullity. Dua ruang isomorf tidak dapat dibedakan
oleh alat bab ini; dimensi menjadi satu-satunya bilangan yang perlu dihitung untuk mengenali
ruang berdimensi hingga.

Ada jebakan tafsir. Isomorfisme tadi bergantung pada pemilihan
basis; tidak ada isomorfisme kanonik antara $V$ dan $\mathbb{R}^{n}$. Karena itu ``$V$ sama
dengan $\mathbb{R}^{n}$'' tidak benar, sedangkan ``$V$ isomorf dengan $\mathbb{R}^{n}$''
benar.
Polinomial berderajat kurang dari tiga dan tripel bilangan real adalah objek berbeda jenis,
tetapi keduanya isomorf sebab setiap polinomial ditentukan tiga koefisiennya.

Sekarang kumpulan semua peta linear antara dua ruang; himpunan itu sendiri ruang vektor.

\begin{teorema}
Himpunan $\mathcal{L}(V,W)$ dari semua peta linear $V\to W$, dengan operasi
$(S+T)(\mathbf{v})=S(\mathbf{v})+T(\mathbf{v})$ dan $(cT)(\mathbf{v})=cT(\mathbf{v})$,
adalah
ruang vektor. Bila $\dim V=n$ dan $\dim W=m$, maka $\dim\mathcal{L}(V,W)=mn$. Pemetaan
$T\mapsto[T]_{C\leftarrow B}$ adalah isomorfisme $\mathcal{L}(V,W)\to M_{m\times n}(\mathbb{R})$.
\end{teorema}

\begin{proof}
$\mathcal{L}(V,W)$ termuat dalam ruang semua fungsi $V\to W$, yang merupakan ruang vektor
menurut contoh di Bab 6. Peta nol linear. Jumlah dua peta linear linear, dan kelipatan
skalar
peta linear juga linear; jadi $\mathcal{L}(V,W)$ subruang, dan karena itu ruang vektor.

Untuk klaim isomorfisme, misalkan $\Phi(T)=[T]_{C\leftarrow B}$. Kelinearan $\Phi$: kolom
ke-$j$ dari $\Phi(S+cT)$ adalah $[S(\mathbf{b}_{j})+cT(\mathbf{b}_{j})]_{C}$, yang sama
dengan
$[S(\mathbf{b}_{j})]_{C}+c[T(\mathbf{b}_{j})]_{C}$ karena koordinat linear. Jadi
$\Phi(S+cT)=\Phi(S)+c\Phi(T)$.

Injektivitas: bila $\Phi(T)$ matriks nol, semua $[T(\mathbf{b}_{j})]_{C}=\mathbf{0}$, jadi
$T(\mathbf{b}_{j})=\mathbf{0}$ untuk semua $j$, dan $T=\mathbf{0}$ karena peta linear
ditentukan oleh nilainya pada basis.

Surjektivitas: beri matriks $A=(a_{ij})$ berukuran $m\times n$. Definisikan
$T(\mathbf{b}_{j})=\sum_{i}a_{ij}\mathbf{c}_{i}$ dan perluas secara linear; menurut
proposisi
tentang nilai pada basis, perluasan itu ada dan tunggal. Kolom ke-$j$ dari $\Phi(T)$ adalah
$(a_{ij})_{i}$, jadi $\Phi(T)=A$.

Karena $\Phi$ isomorfisme dan $\dim M_{m\times n}(\mathbb{R})=mn$, maka
$\dim\mathcal{L}(V,W)=mn$.
\end{proof}

Hitungan itu memberi akibat langsung.

\begin{akibat}
$\dim\mathcal{L}(V,V)=n^{2}$ bila $\dim V=n$. Sebagai contoh,
$\dim\mathcal{L}(P_{2},P_{2})=9$ dan $\dim\mathcal{L}(P_{2},\mathbb{R}^{2})=6$.
\end{akibat}

Hasil terakhir itu menjelaskan mengapa matriks bukan sekadar alat hitung. Ruang operator
dan ruang matriks berdimensi sama dan terhubung oleh isomorfisme yang menghormati komposisi,
sehingga setiap pernyataan tentang operator punya padanan tentang matriks. Yang tidak
diwarisi
adalah kesederhanaan bentuk: operator yang sama dapat rumit pada satu basis dan sederhana
pada
basis lain; mencari basis terbaik adalah program Bab 9.

\begin{tantangan}
Cari satu peta linear $T\colon\mathbb{R}^{3}\to\mathbb{R}^{3}$ yang memetakan persegi satuan
pada bidang $xy$ ke persegi satuan lain dengan luas tetap $1$, tetapi
$\operatorname{rank}(T)=2$. Jelaskan mengapa luas yang bertahan pada satu bidang tidak
memaksa
rank penuh.

Petunjuk 1: periksa satu bidang, bukan seluruh ruang. Setelah itu tanyakan ke mana sumbu
ketiga
pergi.

Petunjuk 2: letakkan rotasi $90^{\circ}$ pada dua koordinat pertama dan nolkan koordinat
ketiga,
sehingga $T(\mathbf{e}_{3})=\mathbf{0}$. Hitung rank dari bayangan $\mathbf{e}_{1}$ dan
$\mathbf{e}_{2}$.

Kalau mentok, kembali ke bagian Transformasi linear, inti, dan jangkauan, khususnya definisi
rank dan akibat rank-nullity.
\end{tantangan}

\begin{untukanak}
Dua kelas di sekolah bisa berbeda ukuran, tetapi punya jumlah kursi yang sama. Kalau
jumlahnya
sama, setiap murid kelas pertama dapat dipasangkan dengan satu kursi di kelas kedua, tanpa
kursi
kosong dan tanpa murid kebagian dua kursi. Untuk urusan menghitung kursi, kedua kelas
setara.

Yang tidak ikut setara adalah nama murid dan wajah mereka. Pemetaan kursi tidak mengatakan
siapa duduk di mana; ia hanya memasangkan. Orang yang berbeda bisa membuat pemetaan berbeda,
dan semuanya sah.

Di bab ini kita akan tahu kapan dua ruang yang tampak berbeda boleh dianggap sama untuk
keperluan hitungan, dan kapan anggapan itu menyesatkan. Kesetaraan itu selalu bergantung
pada
pilihan kita, bukan pada sesuatu yang sudah ada di ruangnya.

Kalau dua kelas punya jumlah kursi sama, apakah kepala sekolah boleh menukar nama kelas
begitu
saja? Jawabannya tidak, dan alasannya ada di bab ini.
\end{untukanak}

\section{Contoh hitung}
\tingkat{menengah}

Tinjau matriks
\[
  A=\begin{pmatrix}1&2&1\\ 2&5&3\\ 3&7&4\end{pmatrix},
  \qquad \text{baris ketiga}=\text{baris pertama}+\text{baris kedua}.
\]
Kita tentukan rank, nullity, dan kedua basisnya, lalu periksa operator yang sama pada basis
lain.\index{rank}\index{nullity}

Eliminasi baris:
\[
  \begin{pmatrix}1&2&1\\ 2&5&3\\ 3&7&4\end{pmatrix}
  \xrightarrow[R_{3}\leftarrow R_{3}-3R_{1}]{R_{2}\leftarrow R_{2}-2R_{1}}
  \begin{pmatrix}1&2&1\\ 0&1&1\\ 0&1&1\end{pmatrix}
  \xrightarrow{R_{3}\leftarrow R_{3}-R_{2}}
  \begin{pmatrix}1&2&1\\ 0&1&1\\ 0&0&0\end{pmatrix}.
\]
Lalu $R_{1}\leftarrow R_{1}-2R_{2}$ memberi bentuk eselon tereduksi
\[
  \begin{pmatrix}1&0&-1\\ 0&1&1\\ 0&0&0\end{pmatrix}.
\]
Ada dua baris tak nol, jadi $\operatorname{rank}A=2$ dan $\det A=0$; kolom pivot adalah kolom
pertama dan kedua.

Inti: dari $x_{1}-x_{3}=0$ dan $x_{2}+x_{3}=0$ kita peroleh $x_{1}=x_{3}$ dan
$x_{2}=-x_{3}$.
Jadi
\[
  \ker A=\operatorname{span}\{(1,-1,1)^{\top}\},\qquad \operatorname{nullity}(A)=1 .
\]
Rank-nullity terpenuhi: $3=2+1$.

Jangkauan: kolom-kolom $A$ adalah $\mathbf{c}_{1}=(1,2,3)^{\top}$,
$\mathbf{c}_{2}=(2,5,7)^{\top}$, dan $\mathbf{c}_{3}=(1,3,4)^{\top}$. Kolom pivot pertama
dan
kedua memberi basis $\{\mathbf{c}_{1},\mathbf{c}_{2}\}$. Kolom ketiga tidak bebas:
$\mathbf{c}_{2}-\mathbf{c}_{1}=\mathbf{c}_{3}$, dan relasi itu persis vektor inti
$(1,-1,1)$: $\mathbf{c}_{1}-\mathbf{c}_{2}+\mathbf{c}_{3}=\mathbf{0}$.

Sekarang pandang $T(\mathbf{x})=A\mathbf{x}$ dan basis
$B'=(\mathbf{b}'_{1},\mathbf{b}'_{2},\mathbf{b}'_{3})$ dengan
$\mathbf{b}'_{1}=(1,1,0)$, $\mathbf{b}'_{2}=(1,0,1)$, $\mathbf{b}'_{3}=(0,1,1)$. Matriks
transisinya dari basis standar $B$ adalah
\[
  P=P_{B\leftarrow B'}=\begin{pmatrix}1&1&0\\ 1&0&1\\ 0&1&1\end{pmatrix},
  \qquad
  P^{-1}=\frac{1}{2}\begin{pmatrix}1&1&-1\\ 1&-1&1\\ -1&1&1\end{pmatrix}.
\]
Bayangan vektor basis baru dihitung lebih dulu:
\[
  A\mathbf{b}'_{1}=\begin{pmatrix}3\\7\\10\end{pmatrix},\quad
  A\mathbf{b}'_{2}=\begin{pmatrix}2\\5\\7\end{pmatrix},\quad
  A\mathbf{b}'_{3}=\begin{pmatrix}3\\8\\11\end{pmatrix}.
\]
Maka
\[
  AP=\begin{pmatrix}3&2&3\\ 7&5&8\\ 10&7&11\end{pmatrix},
  \qquad
  [T]_{B'}=P^{-1}AP=\begin{pmatrix}0&0&0\\ 3&2&3\\ 7&5&8\end{pmatrix}.
\]
Periksa satu entri: baris pertama kolom pertama $\tfrac{1}{2}(3+7-10)=0$; baris kedua
kolom ketiga $\tfrac{1}{2}(3-8+11)=3$. Sisanya mengikuti pola yang sama.

Baris pertama yang seluruhnya nol bukan kebetulan: ketiga bayangan berada di dalam
$\operatorname{span}(\mathbf{b}'_{2},\mathbf{b}'_{3})$, misalnya
$A\mathbf{b}'_{3}=(3,8,11)=3\mathbf{b}'_{2}+8\mathbf{b}'_{3}$. Karena itu koordinat pertama
setiap bayangan nol, dan jangkauan $T$ termuat dalam bidang koordinat pertama.

Pola itu umum. Baris nol menandakan jangkauan yang termuat dalam subruang koordinat
berdimensi lebih rendah; kolom nol menandakan vektor basis yang masuk ke inti. Kedua gejala
hanya tampak pada basis yang dipilih, sedangkan rank dan nullity tidak berubah.

Pemeriksaan menyeluruh:
\[
  \operatorname{rank}[T]_{B'}=2,\qquad \operatorname{trace}[T]_{B'}=0+2+8=10=\operatorname{trace}A,
  \qquad \det[T]_{B'}=0=\det A .
\]
Kesamaan trace bukan kebetulan aritmetika. Entri diagonal ke-$j$ adalah koefisien
$\mathbf{b}'_{j}$ pada bayangan $T(\mathbf{b}'_{j})$, yaitu seberapa besar $T$
mempertahankan
arah itu sendiri. Jumlahnya tidak bergantung pada arah yang dipilih, hanya pada operatornya,
sehingga trace dipakai sebagai sidik jari murah: trace berbeda berarti tidak serupa, trace
sama
belum menjamin apa pun.

Untuk $\mathbf{v}=(3,8,11)$ kita punya $[\mathbf{v}]_{B'}=(0,3,8)^{\top}$ dan
\[
  [T]_{B'}[\mathbf{v}]_{B'}=\begin{pmatrix}0\\30\\79\end{pmatrix},
\]
sedangkan $T(\mathbf{v})=A\mathbf{v}=(30,79,109)^{\top}$ dan koordinatnya
$(0,30,79)^{\top}$. Keduanya cocok.

Peta fitur $G\colon\mathbb{R}^{4}\to\mathbb{R}^{3}$
diberi matriks
\[
  G=\begin{pmatrix}1&2&0&1\\ 0&1&1&1\\ 1&3&1&2\end{pmatrix},
\]
dengan baris ketiga sama dengan jumlah dua baris pertama. Eliminasi memberi bentuk eselon
tereduksi
\[
  \begin{pmatrix}1&0&-2&-1\\ 0&1&1&1\\ 0&0&0&0\end{pmatrix},
\]
jadi $\operatorname{rank}G=2$ dan $\operatorname{nullity}(G)=4-2=2$. Dari
$x_{1}=2x_{2}+x_{4}$ dan $x_{3}=-x_{2}-x_{4}$:
\[
  \mathbf{k}_{1}=(-2,1,-1,0)^{\top},\qquad \mathbf{k}_{2}=(-1,0,-1,1)^{\top}.
\]
Keduanya memenuhi $G\mathbf{k}_{i}=\mathbf{0}$ dan bebas linear, jadi perubahan yang tidak
terlihat oleh ketiga fitur berbentuk bidang dua dimensi, bukan satu garis: penyerang dapat
menggeser masukan ke arah mana pun di bidang itu tanpa menggerakkan skor.

\kertas{Ulangi perhitungan dengan basis $B''=(\mathbf{b}'_{2},\mathbf{b}'_{3},\mathbf{b}'_{1})$,
yaitu $B'$ yang dipermutasi. Ramalkan bagaimana baris nol dan entri lain berpindah, lalu
hitung
$[T]_{B''}$ lewat $Q^{-1}[T]_{B'}Q$ dengan $Q$ matriks permutasi yang sesuai. Periksa bahwa
rank
dan trace tidak berubah, dan jelaskan mengapa $\det A=0$ dapat disimpulkan langsung dari ``baris
ketiga adalah jumlah dua baris pertama''.}

Semua angka di atas diperiksa dengan aritmetika eksak (SymPy); pemeriksaan numpy dan
scipy ada di bagian berikut, termasuk satu angka yang tidak persis nol karena pembulatan
mesin.

\section{Lab komputasi}
\tingkat{lanjut}

\begin{lab}
\begin{verbatim}
import numpy as np
from scipy.linalg import null_space

A = np.array([[1., 2., 1.],
              [2., 5., 3.],
              [3., 7., 4.]])
print("rank A  :", np.linalg.matrix_rank(A))
print("det A   :", np.linalg.det(A))
K = null_space(A)
print("dim ker :", K.shape[1])
print("basis ker:\n", np.round(K, 4))
print("A @ K   :\n", np.round(A @ K, 12))

P = np.array([[1., 1., 0.],
              [1., 0., 1.],
              [0., 1., 1.]])
B = np.linalg.inv(P) @ A @ P
print("[T]_B'  :\n", np.round(B, 10))
print("selisih :", np.linalg.norm(B - np.linalg.inv(P) @ A @ P))
print("trace   :", round(np.trace(B), 10), round(np.trace(A), 10))
print("rank B  :", np.linalg.matrix_rank(B))
print("cond P  :", np.linalg.cond(P))

v = np.array([3., 8., 11.])
print("[v]_B'  :", np.round(np.linalg.solve(P, v), 4))
print("[Av]_B' :", np.round(np.linalg.solve(P, A @ v), 4))
print("B[v]_B' :", np.round(B @ np.linalg.solve(P, v), 4))

x0 = np.array([1., 0., 0.])
d = 0.5 * K[:, 0]
print("A@d     :", np.round(A @ d, 12))
print("geser   :", np.linalg.norm(A @ (x0 + d) - A @ x0))
\end{verbatim}
Keluaran sebenarnya:
\begin{verbatim}
rank A  : 2
det A   : 6.66133814775094e-16
dim ker : 1
basis ker:
 [[ 0.5774]
 [-0.5774]
 [ 0.5774]]
A @ K   :
 [[-0.]
 [ 0.]
 [ 0.]]
[T]_B'  :
 [[0. 0. 0.]
 [3. 2. 3.]
 [7. 5. 8.]]
selisih : 0.0
trace   : 10.0 10.0
rank B  : 2
cond P  : 2.0000000000000004
[v]_B'  : [0. 3. 8.]
[Av]_B' : [ 0. 30. 79.]
B[v]_B' : [ 0. 30. 79.]
A@d     : [-0.  0.  0.]
geser   : 6.280369834735101e-16
\end{verbatim}
\end{lab}

Empat hal perlu dibaca dari keluaran itu.

Pertama, $\det A$ bukan nol melainkan $6{,}66\times10^{-16}$. Matriks ini singular
secara eksak, tetapi aritmetika pecahan biner tidak mengenal nol eksak.
\texttt{matrix\_rank} membandingkan nilai singular terbesar dan terkecil, bukan ambang
sederhana, sehingga rank tetap dilaporkan $2$; Bab 22 membahas mengapa ambang itu perlu.

Pilihan nilai singular, bukan determinan: determinan mengukur volume, dan volume dapat
sangat
kecil tanpa rank berubah. Yang menentukan rank adalah perbandingan nilai singular terkecil
dan
terbesar.

Kedua, \texttt{null\_space} mengembalikan vektor satuan dengan entri $0{,}5774$,
$-0{,}5774$, dan $0{,}5774$, yaitu $(1,-1,1)^{\top}$ dibagi $\sqrt{3}$, dan
$1/\sqrt{3}=0{,}57735$ sesuai. Arahnya sama dengan hasil tangan; panjangnya dinormalkan
karena
scipy memilih basis ortonormal. Basis kernel tidak tunggal, dimensinya tunggal.

Ketiga, $[T]_{B'}$ yang dilaporkan numpy sama persis dengan hasil tangan, dan trace
kedua matriks sama, $10$. Selisih $P^{-1}AP$ yang dihitung dua kali adalah $0{,}0$ karena
$P$
berisi bilangan bulat kecil. Bilangan kondisi $P$ bernilai $2{,}0000000000000004$, cukup
kecil;
basis buruk memperbesar galat, dan itu dibahas di Bab 10.

Keempat, arah inti benar-benar buta bagi $T$. Mendorong $\mathbf{x}_{0}=(1,0,0)$ sejauh
setengah vektor inti mengubah $\mathbf{x}_{0}$ secara nyata, tetapi $A\mathbf{x}$ hanya
bergeser
$6{,}28\times10^{-16}$, yaitu nol dalam presisi mesin.

Peringatan terakhir menyangkut pilihan basis. Basis yang nyaris bergantung membuat koordinat
meledak, dan angka yang dilaporkan kehilangan makna.

\begin{lab}
\begin{verbatim}
Pb = np.array([[1., 1., 0.],
               [1., 1.001, 0.],
               [0., 0., 1.]])
print("cond Pb:", np.linalg.cond(Pb))
print("koord  :", np.linalg.inv(Pb) @ v)
print("galat  :", np.linalg.norm(np.linalg.inv(Pb) @ (Pb @ v) - v))
\end{verbatim}
Keluaran sebenarnya:
\begin{verbatim}
cond Pb: 4002.00075012484
koord  : [-4997.  5000.    11.]
galat  : 4.823324033076557e-13
\end{verbatim}
\end{lab}

Dua vektor basis pertama berbeda hanya pada digit ketiga, dan bilangan kondisi $P_{b}$
mencapai $4002$. Koordinat $(3,8,11)$ meledak menjadi $(-4997,\,5000,\,11)$. Galat
rekonstruksi
masih $4{,}8\times10^{-13}$, tetapi galat relatif pada koordinat membesar sebesar faktor
bilangan
kondisi; dengan 16 digit presisi, sisanya tinggal sekitar 12. Pada basis yang lebih buruk
sisa
itu habis dan rank yang dilaporkan tidak dapat dipercaya.

\begin{untukdosen}
Urutan mengajar yang berhasil dipakai: koordinat lebih dulu, baru peta linear. Mahasiswa
yang
mengenal koordinat sebelum definisi abstrak lebih siap menerima matriks representasi, karena
matriks muncul sebagai daftar koordinat, bukan benda baru. Setelah rank-nullity, masuk
perubahan basis.

Dua miskonsepsi khas. Pertama, menyamakan vektor dengan daftar koordinatnya, lalu bingung
ketika dua basis memberi dua daftar untuk vektor yang sama. Kedua, arah matriks transisi
terbalik. Cara cepat memeriksa: kolom harus berisi alamat vektor basis, dan subskrip dibaca
kanan ke kiri.

Aktivitas papan tulis: beri satu operator pada $\mathbb{R}^{2}$ dan dua basis. Minta separuh
kelas menghitung $[T]_{B'}$ lewat definisi kolom, separuh lewat $P^{-1}AP$, lalu bandingkan
di
papan; kegaduhan yang muncul biasanya tepat pada arah $P$.

Tanda kesiapan lanjut: mahasiswa dapat menjelaskan mengapa sebuah baris nol muncul pada
matriks yang baru, tanpa menghitung ulang. Bila mereka menghubungkannya dengan jangkauan
yang
termuat di subruang koordinat, mereka siap masuk Bab 9.
\end{untukdosen}

\section{Ringkasan}
\tingkat{dasar}

\begin{ringkasan}
\begin{itemize}
  \item Koordinat $[\mathbf{v}]_{B}$ adalah daftar koefisien tunggal terhadap basis terurut;
        peta $\mathbf{v}\mapsto[\mathbf{v}]_{B}$ linear dan bijektif, tetapi bergantung pada
        basis.
  \item Matriks transisi memuat koordinat basis lama dalam basis baru:
        $[\mathbf{v}]_{C}=P_{C\leftarrow B}[\mathbf{v}]_{B}$, dengan
        $P_{B\leftarrow C}=P_{C\leftarrow B}^{-1}$.
  \item Peta linear ditentukan nilainya pada basis; $T$ injektif tepat ketika
        $\ker T=\{\mathbf{0}\}$.
  \item Rank-nullity: $\dim V=\dim\ker T+\dim\operatorname{im}T$; pada dimensi yang sama,
        injektif setara dengan surjektif.
  \item Matriks representasi memenuhi
$[T(\mathbf{v})]_{C}=[T]_{C\leftarrow B}[\mathbf{v}]_{B}$, komposisi menjadi perkalian
        matriks, dan $[T]_{B'}=P^{-1}[T]_{B}P$.
  \item $\dim\mathcal{L}(V,W)=mn$ bila $\dim V=n$, $\dim W=m$; dua ruang berdimensi hingga
        isomorf tepat ketika dimensinya sama.
\end{itemize}
\end{ringkasan}

\section{Latihan}
\tingkat{menengah}

\begin{latihan}
\begin{enumerate}
  \item Dengan $B'=(\mathbf{b}'_{1},\mathbf{b}'_{2},\mathbf{b}'_{3})$ seperti pada contoh
hitung, tentukan $[\mathbf{w}]_{B'}$ untuk $\mathbf{w}=(2,-1,3)$, lalu periksa hasilnya
        dengan mengalikan balik.

\item Buktikan bahwa setiap isomorfisme $V\to\mathbb{R}^{n}$ berbentuk $\varphi_{B}$ untuk
        suatu basis $B$, dan jelaskan mengapa isomorfisme ke $\mathbb{R}^{n}$ tidak tunggal.
        (Petunjuk: ambil prapeta basis standar, lalu tunjukkan himpunan itu basis.)

  \item Untuk $A=\begin{pmatrix}0&1&0\\ 0&0&2\\ 0&0&0\end{pmatrix}$, tentukan rank, nullity,
        basis inti, dan basis jangkauan. Jelaskan hubungannya dengan operator turunan pada
        $P_{2}$.

\item Buktikan $\operatorname{rank}(S\circ T)\le\min(\operatorname{rank}S,\operatorname{rank}T)$
        untuk peta linear yang terdefinisi. Gunakan
        $\operatorname{im}(S\circ T)\subseteq\operatorname{im}S$ dan bandingkan inti.
        (Petunjuk: untuk batas kedua, bandingkan $\ker T$ dengan $\ker(S\circ T)$ lewat
        rank-nullity.)

\item Nilai klaim ini: \emph{dua matriks persegi berukuran sama yang mempunyai rank sama
pasti serupa.} Tentukan benar, salah, atau belum cukup bukti, dengan alasan dan contoh
        penyangkal bila perlu.

\item Misalkan $T\colon\mathbb{R}^{3}\to\mathbb{R}^{2}$ dengan
        $T(x,y,z)=(x+2y,\;y+z)$. Hitung $[T]_{C\leftarrow B}$ untuk
$B=(\,(1,0,0),(1,1,0),(1,1,1)\,)$ dan $C=(\,(1,1),(1,0)\,)$, lalu periksa rumusnya pada
        $\mathbf{v}=(1,2,3)$.
        (Petunjuk: $[T]_{C\leftarrow B}=C^{-1}AB$; hitung $AB$ lebih dulu, lalu selesaikan
        $C\mathbf{x}=\mathbf{y}$ untuk setiap kolomnya.)

  \item Buktikan: bila $U$ subruang $V$ berdimensi hingga dan $\dim U=\dim V$, maka $U=V$.

  \item Buktikan bahwa bila $A$ dan $B$ serupa maka $\operatorname{nullity}(A)
        =\operatorname{nullity}(B)$. Jelaskan mengapa pernyataan itu setara dengan
        ``inti tidak bergantung pada basis''.

  \item Peta fitur $F\colon\mathbb{R}^{4}\to\mathbb{R}^{3}$ diberikan oleh matriks
        \[
          F=\begin{pmatrix}1&0&2&1\\ 0&1&1&3\\ 1&1&3&4\end{pmatrix}.
        \]
        Tentukan rank dan nullity $F$, beri satu vektor inti tak nol, dan jelaskan mengapa
        vektor itu merupakan arah perubahan yang tidak menggeser satu pun fitur.
\end{enumerate}
\end{latihan}

\begin{untukpeneliti}
Yang terbukti di bab ini adalah pernyataan berdimensi hingga atas medan $\mathbb{R}$:
rank-nullity, ketunggalan matriks representasi, kesetaraan isomorfisme dengan kesamaan
dimensi,
dan $\dim\mathcal{L}(V,W)=mn$. Semuanya teorema, bukan konvensi. Klaim ``matriks adalah
transformasinya'' hanya populer; tanpa pemilihan basis klaim itu tidak bermakna, dan contoh
baris
nol di atas menunjukkan bentuk matriks bergantung pada basis.

Sumber primer: Axler, \emph{Linear Algebra Done Right} edisi ke-4, bagian 3A sampai 3D,
memuat
ruang peta linear, inti dan jangkauan, matriks, isomorfisme, serta perubahan basis;
rank-nullity
di sana disebut teorema fundamental untuk peta linear. Hoffman dan Kunze, \emph{Linear
Algebra}
edisi ke-2, Bab 3, lebih rapat namun setara.

Pertanyaan terbuka yang jujur. Klasifikasi operator hingga similaritas tidak diberikan di
sini; Bab 9 hanya sampai diagonalisasi. Secara numerik, keputusan rank bergantung pada
ambang dan
bilangan kondisi; bab ini memakai hasil pustaka tanpa membahas stabilitasnya, utang yang
dibayar
di Bab 22. Batas keberlakuan: di ruang berdimensi tak hingga seluruh argumen runtuh,
rank-nullity
gagal, dan isomorfisme tidak lagi setara dengan kesamaan dimensi.
\end{untukpeneliti}

\peta{Bab ini mengikuti Anton, \emph{Elementary Linear Algebra} edisi ke-11, bagian 4.4, 4.6,
4.8, dan 8.1 sampai 8.5, dengan padanan abstrak di Axler, \emph{Linear Algebra Done Right}
edisi
ke-4, Bab 3, dan Hoffman dan Kunze, \emph{Linear Algebra} edisi ke-2, Bab 3.}